% Modelo de negocio, primera entrega. Máximo 5 páginas, contando el canvas.
% Sustituir cada texto gris (\pendiente{...}) por el contenido del equipo.
% No cambiar los títulos ni el orden de las secciones.
% El canvas se edita en canvas.tex y se inserta aquí con \input. La tabla de
% hipótesis se genera sola a partir de hipotesis.csv: se actualiza el CSV, no
% este archivo. Los resultados de las pruebas van en validacion.tex.
% Compilar dos veces con: pdflatex modelo_negocio.tex
\documentclass[11pt]{article}
\usepackage[utf8]{inputenc}
\usepackage[T1]{fontenc}
\usepackage[spanish,es-noindentfirst,es-nodecimaldot,es-tabla]{babel}
\usepackage[letterpaper,margin=2.2cm,top=1.6cm]{geometry}
\usepackage{xcolor,graphicx,tikz,amssymb}
\usepackage{booktabs,tabularx,array,titling,enumitem,float,pdflscape,afterpage,datatool}
\usepackage[font=small,labelfont=bf,skip=4pt]{caption}
\usepackage[hidelinks]{hyperref}

\definecolor{unamAzul}{HTML}{003D79}
\definecolor{unamOro}{HTML}{D59F0F}
\newcommand{\pendiente}[1]{\textcolor{gray}{\emph{#1}}}
\newcolumntype{L}[1]{>{\raggedright\arraybackslash}p{#1}}
\setlist{nosep,leftmargin=*}
\setlength{\parskip}{4pt}
\setlength{\parindent}{0pt}

% ---------- Hipótesis ----------
% csv-content=literal permite %, &, $ y _ en los textos del CSV. Las claves
% siguen el orden de columnas de la plantilla, que no se cambia.
\DTLread[name=hip,format=csv,csv-content=literal,
  keys={id,bloque,hipotesis,prueba,criterio,estado,evidencia,cambio}]{hipotesis.csv}

% ---------- Encabezado institucional ----------
\pretitle{%
  \begin{center}\small
    Universidad Nacional Autónoma de México\\
    Posgrado en Ciencia e Ingeniería de la Computación\\[2pt]
    Inteligencia Artificial, semestre 2027-1
  \end{center}
  \vspace{-4pt}{\color{unamOro}\rule{\textwidth}{1.2pt}}\par\vspace{14pt}
  \begin{center}{\small\scshape Modelo de negocio}\\[6pt]\LARGE\bfseries\color{unamAzul}}
\posttitle{\par\end{center}\vspace{-4pt}}
\preauthor{\begin{center}\large\begin{tabular}[t]{c}}
\newcommand{\otrorenglon}{\end{tabular}\par\vspace{8pt}\begin{tabular}[t]{c}}
\postauthor{\end{tabular}\par\end{center}}
\predate{\begin{center}\small}
\postdate{\par\end{center}}

\title{GoIA: Navegando la ciudad de forma segura}
\author{%
  Belmont Hernández Héctor Antonio\\ {\small\texttt{habelmont2@gmail.com}}
  \and
  Hernández González Luis Ángel\\ {\small\texttt{luiss.angel92@gmail.com}}
  \otrorenglon
  Núñez Laces Yeni Yuzelmi\\ {\small\texttt{yenilaces@gmail.com}}
  \and
  Villa Aguilar Gerardo Erick\\ {\small\texttt{gerardoerickv@gmail.com}}
}
\date{Versión \pendiente{1}, \today}

\begin{document}
\maketitle
\thispagestyle{empty}

% ---------------------------------------------------------------------------
% EL MODELO EN UN PÁRRAFO. Qué se vende, a quién, qué dolor elimina, cómo llega
% al cliente y cómo se cobra. Quien lea solo este párrafo entiende el negocio.
% ---------------------------------------------------------------------------
\section{El modelo en un párrafo}
\label{sec:resumen}

\pendiente{Un párrafo de cuatro a seis oraciones. Qué vende la empresa, a quién, qué dolor
elimina, por qué canal llega al cliente y cuánto cobra. Cerrar con la hipótesis más riesgosa,
la que el equipo prueba primero.}

La página siguiente muestra el canvas. Cada afirmación del canvas cita la hipótesis que la
respalda, y la Sección~\ref{sec:hipotesis} dice cómo se probará cada hipótesis.

% ---------------------------------------------------------------------------
% CANVAS. Se inserta canvas.tex en una página horizontal, la que sigue a esta.
% El canvas se edita en canvas.tex, no aquí.
% ---------------------------------------------------------------------------
\afterpage{%
  \newgeometry{margin=1cm}%
  \begin{landscape}\thispagestyle{empty}% Business Model Canvas. Una página horizontal dentro de modelo_negocio.tex, que
% lo inserta con % Business Model Canvas. Una página horizontal dentro de modelo_negocio.tex, que
% lo inserta con % Business Model Canvas. Una página horizontal dentro de modelo_negocio.tex, que
% lo inserta con \input{canvas}. Este archivo no se compila solo.
% Llenar los nueve bloques. Cada afirmación termina con el identificador de la
% hipótesis que la respalda en hipotesis.csv, por ejemplo (H1).
% Marcar las hipótesis validadas con \val{H1} y las invalidadas con \inval{H2}.
\definecolor{navy}{HTML}{1F3467}
\providecommand{\val}[1]{\textcolor{green!50!black}{\textbf{(#1 \checkmark)}}}
\providecommand{\inval}[1]{\textcolor{red!70!black}{\textbf{(#1 $\times$)}}}
\providecommand{\pendiente}[1]{\textcolor{gray}{\emph{#1}}}

% ---------- Contenido de los bloques ----------
\newcommand{\Socios}{\begin{itemize}
  \item \pendiente{Proveedores y aliados sin los que el negocio no opera.} (H?)
\end{itemize}}
\newcommand{\Actividades}{\begin{itemize}
  \item \pendiente{Lo que el equipo hace todos los días para entregar valor.} (H?)
\end{itemize}}
\newcommand{\Recursos}{\begin{itemize}
  \item \pendiente{Datos, modelos, personas e infraestructura.} (H?)
\end{itemize}}
\newcommand{\Propuesta}{\begin{itemize}
  \item \pendiente{Qué dolor elimina o qué ganancia crea, para quién.} (H?)
\end{itemize}}
\newcommand{\Relacion}{\begin{itemize}
  \item \pendiente{Autoservicio, asistencia personal, comunidad.} (H?)
\end{itemize}}
\newcommand{\Canales}{\begin{itemize}
  \item \pendiente{Cómo se entera, compra y recibe el producto el cliente.} (H?)
\end{itemize}}
\newcommand{\Segmentos}{\begin{itemize}
  \item \pendiente{Perfil del cliente que paga y del usuario, si son distintos.} (H?)
\end{itemize}}
\newcommand{\Costos}{\begin{itemize}
  \item \pendiente{Costos fijos y variables principales (ver modelo financiero).} (H?)
\end{itemize}}
\newcommand{\Ingresos}{\begin{itemize}
  \item \pendiente{Qué se cobra, cuánto y con qué frecuencia.} (H?)
\end{itemize}}

% ---------- Dibujo. No editar debajo de esta línea ----------
\begingroup
\fontsize{9}{11}\selectfont
\setlist{nosep,leftmargin=1em}
\noindent{\Large\bfseries\color{navy} Business Model Canvas}\hfill
{\large \pendiente{Nombre de la empresa}}\hfill \pendiente{Versión y fecha}

\vspace{4pt}
\noindent
\resizebox{\linewidth}{!}{%
\begin{tikzpicture}[x=1cm,y=1cm]
  \def\W{25.8}\def\H{16.8}\def\h{5.0}
  \tikzset{caja/.style={draw=navy,line width=0.8pt},
           tit/.style={anchor=north west,font=\bfseries\color{navy}},
           cue/.style={anchor=north west,text width=#1cm,align=left}}
  % Retícula: 5 columnas arriba, 2 abajo
  \def\c{5.16}
  \draw[caja] (0,0) rectangle (\W,\H);
  \draw[caja] (0,\h) -- (\W,\h);
  \foreach \i in {1,2,3,4} \draw[caja] (\i*\c,\h) -- (\i*\c,\H);
  \draw[caja] (\W/2,0) -- (\W/2,\h);
  \pgfmathsetmacro\m{(\h+\H)/2}
  \draw[caja] (\c,\m) -- (2*\c,\m);
  \draw[caja] (3*\c,\m) -- (4*\c,\m);
  % Títulos y contenido
  \node[tit] at (0.1,\H-0.1) {Socios clave};
  \node[cue=4.9] at (0.1,\H-0.7) {\Socios};
  \node[tit] at (\c+0.1,\H-0.1) {Actividades clave};
  \node[cue=4.9] at (\c+0.1,\H-0.7) {\Actividades};
  \node[tit] at (\c+0.1,\m-0.1) {Recursos clave};
  \node[cue=4.9] at (\c+0.1,\m-0.7) {\Recursos};
  \node[tit] at (2*\c+0.1,\H-0.1) {Propuesta de valor};
  \node[cue=4.9] at (2*\c+0.1,\H-0.7) {\Propuesta};
  \node[tit] at (3*\c+0.1,\H-0.1) {Relación con clientes};
  \node[cue=4.9] at (3*\c+0.1,\H-0.7) {\Relacion};
  \node[tit] at (3*\c+0.1,\m-0.1) {Canales};
  \node[cue=4.9] at (3*\c+0.1,\m-0.7) {\Canales};
  \node[tit] at (4*\c+0.1,\H-0.1) {Segmentos de clientes};
  \node[cue=4.9] at (4*\c+0.1,\H-0.7) {\Segmentos};
  \node[tit] at (0.1,\h-0.1) {Estructura de costos};
  \node[cue=12.6] at (0.1,\h-0.7) {\Costos};
  \node[tit] at (\W/2+0.1,\h-0.1) {Fuentes de ingresos};
  \node[cue=12.6] at (\W/2+0.1,\h-0.7) {\Ingresos};
\end{tikzpicture}}

\vspace{2pt}
\noindent\scriptsize Basado en Osterwalder y Pigneur, \emph{Business Model Generation} (2010), CC BY-SA 3.0.
\val{H1} validada,\ \inval{H2} invalidada, (H3) pendiente. Evidencia en \texttt{hipotesis.csv} y \texttt{entrevistas.csv}.
\par\endgroup
. Este archivo no se compila solo.
% Llenar los nueve bloques. Cada afirmación termina con el identificador de la
% hipótesis que la respalda en hipotesis.csv, por ejemplo (H1).
% Marcar las hipótesis validadas con \val{H1} y las invalidadas con \inval{H2}.
\definecolor{navy}{HTML}{1F3467}
\providecommand{\val}[1]{\textcolor{green!50!black}{\textbf{(#1 \checkmark)}}}
\providecommand{\inval}[1]{\textcolor{red!70!black}{\textbf{(#1 $\times$)}}}
\providecommand{\pendiente}[1]{\textcolor{gray}{\emph{#1}}}

% ---------- Contenido de los bloques ----------
\newcommand{\Socios}{\begin{itemize}
  \item \pendiente{Proveedores y aliados sin los que el negocio no opera.} (H?)
\end{itemize}}
\newcommand{\Actividades}{\begin{itemize}
  \item \pendiente{Lo que el equipo hace todos los días para entregar valor.} (H?)
\end{itemize}}
\newcommand{\Recursos}{\begin{itemize}
  \item \pendiente{Datos, modelos, personas e infraestructura.} (H?)
\end{itemize}}
\newcommand{\Propuesta}{\begin{itemize}
  \item \pendiente{Qué dolor elimina o qué ganancia crea, para quién.} (H?)
\end{itemize}}
\newcommand{\Relacion}{\begin{itemize}
  \item \pendiente{Autoservicio, asistencia personal, comunidad.} (H?)
\end{itemize}}
\newcommand{\Canales}{\begin{itemize}
  \item \pendiente{Cómo se entera, compra y recibe el producto el cliente.} (H?)
\end{itemize}}
\newcommand{\Segmentos}{\begin{itemize}
  \item \pendiente{Perfil del cliente que paga y del usuario, si son distintos.} (H?)
\end{itemize}}
\newcommand{\Costos}{\begin{itemize}
  \item \pendiente{Costos fijos y variables principales (ver modelo financiero).} (H?)
\end{itemize}}
\newcommand{\Ingresos}{\begin{itemize}
  \item \pendiente{Qué se cobra, cuánto y con qué frecuencia.} (H?)
\end{itemize}}

% ---------- Dibujo. No editar debajo de esta línea ----------
\begingroup
\fontsize{9}{11}\selectfont
\setlist{nosep,leftmargin=1em}
\noindent{\Large\bfseries\color{navy} Business Model Canvas}\hfill
{\large \pendiente{Nombre de la empresa}}\hfill \pendiente{Versión y fecha}

\vspace{4pt}
\noindent
\resizebox{\linewidth}{!}{%
\begin{tikzpicture}[x=1cm,y=1cm]
  \def\W{25.8}\def\H{16.8}\def\h{5.0}
  \tikzset{caja/.style={draw=navy,line width=0.8pt},
           tit/.style={anchor=north west,font=\bfseries\color{navy}},
           cue/.style={anchor=north west,text width=#1cm,align=left}}
  % Retícula: 5 columnas arriba, 2 abajo
  \def\c{5.16}
  \draw[caja] (0,0) rectangle (\W,\H);
  \draw[caja] (0,\h) -- (\W,\h);
  \foreach \i in {1,2,3,4} \draw[caja] (\i*\c,\h) -- (\i*\c,\H);
  \draw[caja] (\W/2,0) -- (\W/2,\h);
  \pgfmathsetmacro\m{(\h+\H)/2}
  \draw[caja] (\c,\m) -- (2*\c,\m);
  \draw[caja] (3*\c,\m) -- (4*\c,\m);
  % Títulos y contenido
  \node[tit] at (0.1,\H-0.1) {Socios clave};
  \node[cue=4.9] at (0.1,\H-0.7) {\Socios};
  \node[tit] at (\c+0.1,\H-0.1) {Actividades clave};
  \node[cue=4.9] at (\c+0.1,\H-0.7) {\Actividades};
  \node[tit] at (\c+0.1,\m-0.1) {Recursos clave};
  \node[cue=4.9] at (\c+0.1,\m-0.7) {\Recursos};
  \node[tit] at (2*\c+0.1,\H-0.1) {Propuesta de valor};
  \node[cue=4.9] at (2*\c+0.1,\H-0.7) {\Propuesta};
  \node[tit] at (3*\c+0.1,\H-0.1) {Relación con clientes};
  \node[cue=4.9] at (3*\c+0.1,\H-0.7) {\Relacion};
  \node[tit] at (3*\c+0.1,\m-0.1) {Canales};
  \node[cue=4.9] at (3*\c+0.1,\m-0.7) {\Canales};
  \node[tit] at (4*\c+0.1,\H-0.1) {Segmentos de clientes};
  \node[cue=4.9] at (4*\c+0.1,\H-0.7) {\Segmentos};
  \node[tit] at (0.1,\h-0.1) {Estructura de costos};
  \node[cue=12.6] at (0.1,\h-0.7) {\Costos};
  \node[tit] at (\W/2+0.1,\h-0.1) {Fuentes de ingresos};
  \node[cue=12.6] at (\W/2+0.1,\h-0.7) {\Ingresos};
\end{tikzpicture}}

\vspace{2pt}
\noindent\scriptsize Basado en Osterwalder y Pigneur, \emph{Business Model Generation} (2010), CC BY-SA 3.0.
\val{H1} validada,\ \inval{H2} invalidada, (H3) pendiente. Evidencia en \texttt{hipotesis.csv} y \texttt{entrevistas.csv}.
\par\endgroup
. Este archivo no se compila solo.
% Llenar los nueve bloques. Cada afirmación termina con el identificador de la
% hipótesis que la respalda en hipotesis.csv, por ejemplo (H1).
% Marcar las hipótesis validadas con \val{H1} y las invalidadas con \inval{H2}.
\definecolor{navy}{HTML}{1F3467}
\providecommand{\val}[1]{\textcolor{green!50!black}{\textbf{(#1 \checkmark)}}}
\providecommand{\inval}[1]{\textcolor{red!70!black}{\textbf{(#1 $\times$)}}}
\providecommand{\pendiente}[1]{\textcolor{gray}{\emph{#1}}}

% ---------- Contenido de los bloques ----------
\newcommand{\Socios}{\begin{itemize}
  \item \pendiente{Proveedores y aliados sin los que el negocio no opera.} (H?)
\end{itemize}}
\newcommand{\Actividades}{\begin{itemize}
  \item \pendiente{Lo que el equipo hace todos los días para entregar valor.} (H?)
\end{itemize}}
\newcommand{\Recursos}{\begin{itemize}
  \item \pendiente{Datos, modelos, personas e infraestructura.} (H?)
\end{itemize}}
\newcommand{\Propuesta}{\begin{itemize}
  \item \pendiente{Qué dolor elimina o qué ganancia crea, para quién.} (H?)
\end{itemize}}
\newcommand{\Relacion}{\begin{itemize}
  \item \pendiente{Autoservicio, asistencia personal, comunidad.} (H?)
\end{itemize}}
\newcommand{\Canales}{\begin{itemize}
  \item \pendiente{Cómo se entera, compra y recibe el producto el cliente.} (H?)
\end{itemize}}
\newcommand{\Segmentos}{\begin{itemize}
  \item \pendiente{Perfil del cliente que paga y del usuario, si son distintos.} (H?)
\end{itemize}}
\newcommand{\Costos}{\begin{itemize}
  \item \pendiente{Costos fijos y variables principales (ver modelo financiero).} (H?)
\end{itemize}}
\newcommand{\Ingresos}{\begin{itemize}
  \item \pendiente{Qué se cobra, cuánto y con qué frecuencia.} (H?)
\end{itemize}}

% ---------- Dibujo. No editar debajo de esta línea ----------
\begingroup
\fontsize{9}{11}\selectfont
\setlist{nosep,leftmargin=1em}
\noindent{\Large\bfseries\color{navy} Business Model Canvas}\hfill
{\large \pendiente{Nombre de la empresa}}\hfill \pendiente{Versión y fecha}

\vspace{4pt}
\noindent
\resizebox{\linewidth}{!}{%
\begin{tikzpicture}[x=1cm,y=1cm]
  \def\W{25.8}\def\H{16.8}\def\h{5.0}
  \tikzset{caja/.style={draw=navy,line width=0.8pt},
           tit/.style={anchor=north west,font=\bfseries\color{navy}},
           cue/.style={anchor=north west,text width=#1cm,align=left}}
  % Retícula: 5 columnas arriba, 2 abajo
  \def\c{5.16}
  \draw[caja] (0,0) rectangle (\W,\H);
  \draw[caja] (0,\h) -- (\W,\h);
  \foreach \i in {1,2,3,4} \draw[caja] (\i*\c,\h) -- (\i*\c,\H);
  \draw[caja] (\W/2,0) -- (\W/2,\h);
  \pgfmathsetmacro\m{(\h+\H)/2}
  \draw[caja] (\c,\m) -- (2*\c,\m);
  \draw[caja] (3*\c,\m) -- (4*\c,\m);
  % Títulos y contenido
  \node[tit] at (0.1,\H-0.1) {Socios clave};
  \node[cue=4.9] at (0.1,\H-0.7) {\Socios};
  \node[tit] at (\c+0.1,\H-0.1) {Actividades clave};
  \node[cue=4.9] at (\c+0.1,\H-0.7) {\Actividades};
  \node[tit] at (\c+0.1,\m-0.1) {Recursos clave};
  \node[cue=4.9] at (\c+0.1,\m-0.7) {\Recursos};
  \node[tit] at (2*\c+0.1,\H-0.1) {Propuesta de valor};
  \node[cue=4.9] at (2*\c+0.1,\H-0.7) {\Propuesta};
  \node[tit] at (3*\c+0.1,\H-0.1) {Relación con clientes};
  \node[cue=4.9] at (3*\c+0.1,\H-0.7) {\Relacion};
  \node[tit] at (3*\c+0.1,\m-0.1) {Canales};
  \node[cue=4.9] at (3*\c+0.1,\m-0.7) {\Canales};
  \node[tit] at (4*\c+0.1,\H-0.1) {Segmentos de clientes};
  \node[cue=4.9] at (4*\c+0.1,\H-0.7) {\Segmentos};
  \node[tit] at (0.1,\h-0.1) {Estructura de costos};
  \node[cue=12.6] at (0.1,\h-0.7) {\Costos};
  \node[tit] at (\W/2+0.1,\h-0.1) {Fuentes de ingresos};
  \node[cue=12.6] at (\W/2+0.1,\h-0.7) {\Ingresos};
\end{tikzpicture}}

\vspace{2pt}
\noindent\scriptsize Basado en Osterwalder y Pigneur, \emph{Business Model Generation} (2010), CC BY-SA 3.0.
\val{H1} validada,\ \inval{H2} invalidada, (H3) pendiente. Evidencia en \texttt{hipotesis.csv} y \texttt{entrevistas.csv}.
\par\endgroup
\end{landscape}%
  \restoregeometry}

% ---------------------------------------------------------------------------
% CLIENTE Y PROPUESTA DE VALOR. La persona es un personaje ficticio construido
% con hechos de clientes reales, no con suposiciones. Cada celda cita las
% entrevistas (E01) o las hipótesis (H1) que la respaldan. Si quien paga y quien
% usa son la misma persona, borrar la columna Usuario.
% ---------------------------------------------------------------------------
\section{Cliente y propuesta de valor}
\label{sec:cliente}

\pendiente{Una o dos oraciones que presentan el segmento principal y por qué se eligió sobre
otros segmentos posibles.} La Tabla~\ref{tab:persona} describe a la persona que paga y a la que
usa el producto.

\begin{table}[H]
  \caption{Perfil de la persona que paga y de la que usa el producto. Cada celda cita la
    entrevista o la hipótesis que la respalda.}
  \label{tab:persona}
  \small
  \renewcommand{\arraystretch}{1.25}
  \begin{tabularx}{\textwidth}{@{}>{\bfseries\raggedright}p{0.17\textwidth}
                                  >{\raggedright\arraybackslash}X>{\raggedright\arraybackslash}X@{}}
    \toprule
    & Cliente que paga & Usuario \\
    \midrule
    Nombre y cargo & \pendiente{Nombre ficticio y cargo, por ejemplo ``Laura, dueña de la
                     paquetería''.} & \pendiente{Por ejemplo, ``Jorge, jefe de despacho''.} \\
    Organización & \pendiente{Tipo, tamaño y ciudad.} & \\
    Un día típico & \pendiente{Qué hace y en qué momento aparece el problema.} & \\
    Objetivos & \pendiente{Qué intenta lograr en su trabajo.} & \\
    Dolores & \pendiente{Qué le pasa, con qué frecuencia y cuánto le cuesta en tiempo o
              dinero.} (E?) & \\
    Ganancias & \pendiente{Qué resultado lo haría quedar bien o ahorrar.} & \\
    Solución actual & \pendiente{Qué hace hoy para resolverlo y cuánto gasta.} & \\
    Cómo compra & \pendiente{Quién decide, quién aprueba y con qué presupuesto.} & \\
    Dónde se informa & \pendiente{Personas, grupos, eventos y medios en que confía. Estos
                       lugares son candidatos a canal.} & \\
    \bottomrule
  \end{tabularx}
\end{table}

La propuesta de valor es \pendiente{una oración con la forma: para [cliente], que sufre
[dolor], [producto] logra [resultado medible], a diferencia de [solución actual].} (H?)

% ---------------------------------------------------------------------------
% RELACIÓN CON CLIENTES Y CANALES. Los canales cubren las cinco fases del
% recorrido del cliente. Un canal concreto dice a quién llega y con qué mensaje:
% "redes sociales" no basta, "anuncios en LinkedIn a jefes de despacho" sí. La
% relación dice cómo se consigue, se conserva y se hace crecer a cada cliente.
% El costo de cada canal y de cada relación alimenta el CAC del modelo financiero.
% ---------------------------------------------------------------------------
\section{Relación con clientes y canales}
\label{sec:canales}

\pendiente{Una o dos oraciones que dicen cuál es el canal principal para conseguir clientes y
por qué llega adonde el segmento ya está.} La Tabla~\ref{tab:canales} describe los canales en
cada fase, y la Tabla~\ref{tab:relacion} describe la relación con el cliente.

\begin{table}[H]
  \caption{Canales en las cinco fases del recorrido del cliente. Directo significa que el
    equipo opera el canal, e indirecto que lo opera un socio.}
  \label{tab:canales}
  \small
  \renewcommand{\arraystretch}{1.2}
  \begin{tabularx}{\textwidth}{@{}>{\bfseries}p{0.15\textwidth}Xp{0.11\textwidth}p{0.16\textwidth}c@{}}
    \toprule
    Fase & Canal & Tipo & Costo estimado & Hipótesis \\
    \midrule
    Conocimiento & \pendiente{Cómo se entera el cliente de que existe el producto.} & & & \\
    Evaluación & \pendiente{Cómo compara y prueba antes de comprar.} & & & \\
    Compra & \pendiente{Dónde y cómo paga.} & & & \\
    Entrega & \pendiente{Cómo recibe el producto.} & & & \\
    Posventa & \pendiente{Cómo recibe soporte.} & & & \\
    \bottomrule
  \end{tabularx}
\end{table}

\begin{table}[H]
  \caption{Relación con el cliente en tres etapas. Los tipos son asistencia personal,
    asistencia dedicada, autoservicio, servicio automatizado, comunidad y cocreación.}
  \label{tab:relacion}
  \small
  \renewcommand{\arraystretch}{1.2}
  \begin{tabularx}{\textwidth}{@{}>{\bfseries}p{0.13\textwidth}p{0.2\textwidth}Xc@{}}
    \toprule
    Etapa & Tipo & Cómo se hace y cuánto cuesta & Hipótesis \\
    \midrule
    Conseguir & & \pendiente{Por ejemplo, demostración en persona y piloto de tres días.} & \\
    Conservar & & \pendiente{Por ejemplo, soporte por WhatsApp en horario laboral.} & \\
    Crecer & & \pendiente{Por ejemplo, más usuarios por cliente o un segundo producto.} & \\
    \bottomrule
  \end{tabularx}
\end{table}

% ---------------------------------------------------------------------------
% COMPETENCIA. Todo problema real tiene una solución actual, aunque sea una hoja
% de cálculo. Incluir esa solución y al menos dos competidores directos. Los dos
% atributos son los que más le importan al cliente según las entrevistas.
% ---------------------------------------------------------------------------
\section{Competencia}
\label{sec:competencia}

La Tabla~\ref{tab:competencia} compara las opciones del cliente en
\pendiente{atributo 1} y \pendiente{atributo 2}, los dos atributos que más mencionaron las
personas entrevistadas.

\begin{table}[H]
  \caption{Opciones actuales del cliente. Calificación de 1 (bajo) a 5 (alto) en los dos
    atributos que más le importan, y precio con su fuente.}
  \label{tab:competencia}
  \small
  \renewcommand{\arraystretch}{1.2}
  \begin{tabularx}{\textwidth}{@{}Xp{0.17\textwidth}ccp{0.2\textwidth}@{}}
    \toprule
    Opción & Tipo & \pendiente{Atributo 1} & \pendiente{Atributo 2} & Precio y fuente \\
    \midrule
    \pendiente{Solución actual del cliente} & Solución actual & & & \\
    \pendiente{Competidor 1} & Directo & & & \\
    \pendiente{Competidor 2} & Indirecto & & & \\
    \textbf{\pendiente{Nuestra empresa}} & & & & \\
    \bottomrule
  \end{tabularx}
\end{table}

\pendiente{Un párrafo que dice qué sostiene la ventaja frente a la mejor alternativa: datos
propios, un resultado medido o un canal que la competencia no tiene. Una ventaja sin evidencia
es una hipótesis y se registra como tal.}

% ---------------------------------------------------------------------------
% TAMAÑO DE MERCADO. De abajo hacia arriba: número de clientes por precio anual.
% Cada número lleva su fuente. Las cifras son las mismas de la hoja Supuestos del
% modelo financiero.
% ---------------------------------------------------------------------------
\section{Tamaño de mercado}
\label{sec:mercado}

La Tabla~\ref{tab:mercado} estima el mercado de abajo hacia arriba, con el número de clientes
por el precio anual. \pendiente{Una oración que justifica el porcentaje que lleva del TAM al
SAM y del SAM al SOM.}

\begin{table}[H]
  \caption{Tamaño de mercado. TAM es todo el mercado si cada cliente posible comprara, SAM es la
    parte que alcanzan los canales del equipo y SOM es la parte que el equipo espera ganar en 24
    meses. Montos en pesos por año.}
  \label{tab:mercado}
  \small
  \renewcommand{\arraystretch}{1.2}
  \begin{tabularx}{\textwidth}{@{}lrrrX@{}}
    \toprule
    Nivel & Clientes & Precio anual & Mercado & Fuente o supuesto \\
    \midrule
    TAM & & & & \pendiente{Por ejemplo, número de consultorios según el DENUE de INEGI.} \\
    SAM & & & & \pendiente{Por ejemplo, el porcentaje en la CDMX.} \\
    SOM & & & & \pendiente{Por ejemplo, la tasa de conversión medida en la landing page.} \\
    \bottomrule
  \end{tabularx}
\end{table}

% ---------------------------------------------------------------------------
% HIPÓTESIS Y PRUEBAS. La tabla se genera a partir de hipotesis.csv. Cada
% afirmación del canvas tiene una hipótesis, y cada hipótesis tiene una prueba y
% un criterio de éxito fijado antes de probarla. Los resultados van en
% validacion.tex, en la entrega final.
% ---------------------------------------------------------------------------
\section{Hipótesis y pruebas}
\label{sec:hipotesis}

\pendiente{Un párrafo que dice cuáles son las dos o tres hipótesis más riesgosas, las que
harían fracasar el negocio si fueran falsas, y en qué orden se probarán.} La
Tabla~\ref{tab:hipotesis} lista las hipótesis del canvas.

\begin{table}[H]
  \caption{Hipótesis del canvas, con la prueba que la pondrá a prueba y el criterio de éxito
    fijado antes de probarla.}
  \label{tab:hipotesis}
  \footnotesize
  \renewcommand{\arraystretch}{1.2}
  \begin{tabular}{@{}lL{0.38\textwidth}L{0.22\textwidth}L{0.28\textwidth}@{}}
    \toprule
    id & Hipótesis (bloque) & Prueba & Criterio de éxito \\
    \midrule
    \DTLforeach*{hip}{\hId=id,\hBloque=bloque,\hTexto=hipotesis,\hPrueba=prueba,\hCrit=criterio}{%
      \DTLiffirstrow{}{\\}%
      \hId & \hTexto{} {\color{gray}(\hBloque)} & \hPrueba & \hCrit}\\
    \bottomrule
  \end{tabular}
\end{table}

\end{document}
